% ============================================================
% paper/sections/05_thresholds_structural_tension.tex
% Thresholds, f_k Components, and Structural Tension (finalized)
% ============================================================

\section{Thresholds, Components, and Structural Tension}

\subsection{Role of Thresholds in the Continuum}

Thresholds $\Theta_{\mathrm{EA}}$ define the admissibility limits of the
enterprise continuum.
They separate regions of the state space that are structurally
admissible from those in which the continuum is undefined, unstable, or
non-diagnosable.

Thresholds are not targets, objectives, or performance benchmarks.
They are constraints that determine whether the continuum can exist,
persist, and remain diagnosable under the given structure.

\subsection{Structural Tension and Inequality-Based Formulation}

Structural tension arises when internal dynamics drive the enterprise
state toward the boundary $\partial\Omega_{\mathrm{EA}}$ defined by the
thresholds.

In \texttt{ETS.K\_EA.v1.1}, structural tension is expressed through a set
of component functions $f_k$, each evaluated as an inequality.
The formalism operates on the sign, ordering, and logical combination of
these components rather than on calibrated numerical magnitudes.

No numerical values are fixed at the theory level.
Any quantitative instantiation occurs outside the theory and does not
alter the formal structure.

\subsection{Threshold Components}

The Executable Theory Specification defines the following threshold
components:

\begin{itemize}
  \item $f_{\text{exist}}$: existence admissibility,
  \item $f_{\text{stab}}$: stability and persistence,
  \item $f_{\text{inertia}}$: structural inertia,
  \item $f_{\text{collapse}}$: irreversible structural breakdown,
  \item $f_{\text{stop}}$: termination of diagnosis.
\end{itemize}

Each component corresponds to a subset of $\Theta_{\mathrm{EA}}$.
Components are evaluated independently, subject to explicit precedence
rules encoded in the phase logic.

\subsection{Interpretation of Components}

\paragraph{Existence Component ($f_{\text{exist}}$).}
Evaluates whether the admissible state space $\Omega_{\mathrm{EA}}$
remains non-empty.
Violation indicates loss of definability of the enterprise continuum
itself.

\paragraph{Stability Component ($f_{\text{stab}}$).}
Evaluates the integrity of required recurrent cycles over the internal
time scale $\tau_{\mathrm{EA}}$.
Violation indicates loss of persistence even if instantaneous
admissibility holds.

\paragraph{Inertia Component ($f_{\text{inertia}}$).}
Evaluates the ability of the continuum to transition between admissible
states under endogenous dynamics.
Violation indicates structural immobilization without immediate loss of
existence.

\paragraph{Collapse Component ($f_{\text{collapse}}$).}
Signals irreversible loss of structural coherence.
It typically follows sustained violation of stability conditions and
cannot be interpreted as a reversible state.

\paragraph{Termination Component ($f_{\text{stop}}$).}
Encodes conditions under which diagnosis must terminate regardless of
other component values.
This includes loss of observability and violation of compatibility with
parent continua.

\subsection{Absence of Implicit Optimization}

The presence of thresholds does not imply that the enterprise should be
kept distant from them, nor that movement within the admissible region
is desirable or undesirable.

Thresholds define only the boundary of admissible diagnosis.
Crossing certain thresholds renders further analysis undefined within
the model.

\subsection{Structural Tension Versus Performance Stress}

Structural tension is distinct from operational, financial, or
performance-related stress.
An enterprise may experience severe stress while remaining structurally
admissible and diagnosable.

Conversely, structural thresholds may be violated without immediate or
visible performance degradation.
The threshold formalism makes this distinction explicit by separating
structural admissibility from observed outcomes.
